% /chapters/0-preamble/2-publications.tex

\begin{publications}
\begin{itemize}
    \item \textbf{Section~\ref{sec:concise_L0}: A concise proof of the $L_0$ dichotomy}
    \\ \textit{Presentation:} Fields Institute Set Theory Seminar, Toronto (March 28, 2025).
    \\ \textit{Note:} Contains original proofs first disseminated at the Fields Institute.

    \item \textbf{Section~\ref{sec:dichromatic}: No countable basis for Borel directed graphs of dichromatic number at least three}
    \\ \textit{Presentation:} Caltech Logic Seminar (April 29, 2026).
    \\ \textit{Note:} This new result completes the picture for finite Borel chromatic and dichromatic thresholds.

    \item \textbf{Section~\ref{sec:cantor_game}: Uncountable sets and an infinite linear order
game}
    \\ \textit{Source:} Adapted from \cite{MatosWiederholdSalvetti2025}.
    \\ \textit{Presentation:} 2023 Summer Meeting of the Canadian Mathematical Society, Ottawa.
    \\ \textit{Note:} This work provides a solution to a problem originally posed by Brian and Clontz about games.

    \item \textbf{Chapter~\ref{ch:deep_computations}: Deep Computations}
    \\ \textit{Source:} Primarily based on \cite{Duenezetal2026}.
    \\ \textit{Presentation:} A multi-part series at the Fields Institute Set Theory Seminar (2024--2026).
    \\ \textit{Note:} This chapter focuses on the author's contributions to the separation of learnability classes with computational structures and their connections to other fields.

    \item \textbf{Section~\ref{sec:border_tracing}: Border tracing in oriented adjacency graphs of polygonal tilings}
    \\ \textit{Source:} Based on my contributions to \cite{WiederholdMatosWiederhold2026}.

    \item \textbf{Section~\ref{sec:recursive_constructions}: Recursive constructions}
    \\ \textit{Source:} Extends joint work found in \cite{Eslavaetal2025} and \cite{Caroetal2023}.
    \\ \textit{Presentation:} 2022 Summer Meeting of the Canadian Mathematical Society, St John's.
    \\ \textit{Note:} This section presents original generalizations and recursive constructions that extend beyond the currently published literature.

    \item \textbf{Section~\ref{sec:amoeba_trees}: Amoeba trees}
    \\ \textit{Source:} Based on \cite{GomezetalPreprint}.
\end{itemize}
\end{publications}